% OWNER: Lynn (L6). Target: ~0.5 page. Test numbers: README section 5b / results/final/.
% Validation-only observations are labelled as such. Update the TODOs when the ablations (A6) and
% the control baselines (F5) are in.
\section{Discussion}
\label{sec:discussion}

\paragraph{Answer to the research question.}
LLM-guided search over LLM-predicted states does not improve Next Action accuracy over direct
generation under a matched call budget (pooled 95\% intervals: 8B $[-0.13, +0.003]$, 70B $[-0.09, +0.06]$). How the calls are
spent matters: with the same 14 calls on 8B, best-of-$N$ is higher than search in every seed (significantly in one; 70B: tied), and also under a token cap.

\paragraph{Evaluation, more than generation, is the bottleneck.}
The search does find correct actions: one is among the root candidates in 31\% of the questions for
8B and 59\% for 70B, against final accuracies of 0.16 and 0.36. Selection loses much of this potential. 8B ranks a question's candidates at chance (within-question AUC $0.48$), so its search ends below its own first proposal ($0.159$ vs.\ $0.179$). Scale helps selection too: 70B lifts its first proposal from $0.279$ to $0.362$ and prunes 71\% of inapplicable candidates (8B: 14\%), though also 23\% of applicable ones, and picks the correct candidate in 61\% of the questions where one is proposed (a random pick: 43\%; after pruning, 56\%, and the proposer's first remaining candidate 60\%, so pruning, not scoring, does most of this). Its first proposal is weaker than greedy two-shot direct generation ($0.377$), and selection recovers only part of the gap. The 8B ablations point the same way, though only $k{=}5$ differs significantly from the frozen search, in one seed ($p = 0.021$). Trusting the proposer's own order gives $0.190$, above every evaluator ($0.108$--$0.159$) and random scores ($0.133$). Proposing more candidates hurts in every seed ($k{=}5$: $0.115$), plausibly because the evaluator then has more ways to choose wrongly, and removing
state prediction leaves accuracy unchanged at under a third of the calls. Depth adds nothing either: doubling the budget to let other candidates' subtrees compete does not help ($0.154$; Section~\ref{sec:method}). 

\paragraph{Implications.}
Without a verifier, search can only be as good as the model's judgment of its own candidates, and for open-ended planning this judgment is the weak link, as found for self-critique of plans~\citep{valmeekam2023selfcritique,stechly2024selfverification}. The large gap between the candidates found
and the answers chosen suggests that the same search could pay off with a reliable external
evaluator, such as a learned verifier or a cheap symbolic check, as in LLM-Modulo frameworks~\citep{kambhampati2024llmmodulo}, Thought of Search~\citep{katz2024tos} and LLM+P~\citep{liu2023llmp}.
